%!TEX root = ./../../springer-mono-format/book.tex
%!TEX root = ./../../thesis-trier.tex

\label{def:power_aut} \IN{power (semi-)automaton}
Let $\semiautomat{A}{Q}{\Sigma}{\delta}$. The \emph{power (semi-)automaton} 
is $\semiautomat{P({\mathcal A})}{\mathcal P(Q)}{\Sigma}{\delta}$
where $\delta \colon \mathcal P(Q)\times \Sigma \to \Sigma$ is the extension of
$\delta \colon Q \times \Sigma \to Q$ to subsets $\applytrans{\delta}{S}{x} = \{ \applytrans{\delta}{q}{x} \mid q \in Q \}$ 
for $S \subseteq Q$ and $x \in \Sigma$ (see also Section~\ref{intro:sec:basic}).
\Cref{fig:thm_not_usable} shows the power automaton
of a four-state automaton. Also, part of the power 
automaton of the \v{C}ern\'y automaton $\mathcal C_4$
is depicted in \cref{sync_max:fig:C4-poweraut}.

In this section, we present properties of powers of a single word in the power automaton. These can also be interpreted as properties of transformations on subsets of a finite set; however, since our focus is on automata, we express them in automata-theoretic terms.

\todo[inline]{bild kompakter}

\begin{figure}[htb]
\begin{minipage}{0.25\textwidth}
 \scalebox{1}{\begin{tikzpicture}[shorten >=1pt,->,node distance=2.3cm]
  %\tikzstyle{vertex}=[circle,color=red!60,fill=black!10,minimum size=25pt,inner sep=1pt,very thick]
  %\tikzstyle{vertex}=[circle,draw,minimum size=25pt,inner sep=1pt]
  
  \node[state] (0) {$0$};
  \node[state, right of=0] (1) {$1$};
  \node[state, below of=1] (2) {$2$};
  \node[state, left of=2]  (3) {$3$};
  
  \path[->] (0) edge [above] node {$a,b$} (1)
            (1) edge [bend left, right] node {$a,b$} (2)
            (2) edge [below] node {$b$} (3)
            (3) edge [left]  node {$b$} (0)
            (2) edge [bend left, left] node {$a$} (1)
            (3) edge [loop below] node {$a$} (3);
\end{tikzpicture}}%
\end{minipage}%\hspace*{0.25cm}
\begin{minipage}{0.75\textwidth}
\scalebox{1}{\begin{tikzpicture}[shorten >=1pt,->,node distance=1cm and 0.9cm]
  \tikzstyle{vertex}=[]
  
  % \node at (0,-1) (s) {start};  
   
  \node[vertex] (0123)               {$\{0,1,2,3\}$};
  \node[vertex, right = of 0123, xshift=-0.5cm] (123) {$\{1,2,3\}$};
  \node[vertex, below = of 123]  (023) {$\{0,2,3\}$};
  \node[vertex, below = of 023]   (013) {$\{0,1,3\}$};
  \node[vertex, below left = of 013,xshift=0.7cm]  (012) {$\{0,1,2\}$};
  
  \node[vertex, right = of 023] (13) {$\{1,3\}$};
  \node[vertex, above = of 13]  (02) {$\{0,2\}$};
  \node[vertex, below = of 13]  (23) {$\{2,3\}$};
  \node[vertex, right = of 23]  (03) {$\{0,3\}$};
  \node[vertex, below = of 03]  (01) {$\{0,1\}$};
  \node[vertex, left = of 01]  (12) {$\{1,2\}$};
  
  \node[vertex,  right = of 02, yshift=0.2cm] (0) {$\{0\}$};
  \node[vertex,  right = of 0]  (1) {$\{1\}$};
  \node[vertex,  below = of 1]  (2) {$\{2\}$};
  \node[vertex,  left = of 2]   (3) {$\{3\}$};
  
  
  \path (123) edge [loop above] node {$a$} (123);
  \path (012) edge [above] node {$a$} (12);
  \path (013) edge [bend left=45,left] node {$a$} (123);
  
  %\path[->] (s) edge (0123);
  
    \path[->] (0) edge [below] node {$a,b$} (1)
            (1) edge [bend left, right] node {$a,b$} (2)
            (2) edge [below] node {$b$} (3)
            (3) edge [left]  node {$b$} (0)
            (2) edge [bend left, right] node {$a$} (1)
            (3) edge [loop below] node {$a$} (3);
            
  
  \path[->] (0123) edge [above] node {$a$} (123);
  \path[->] (0123) edge [loop above] node {$b$} (0123);
  \path[->] (123)  edge [right] node {$b$} (023); 
  \path[->] (023)  edge [right] node {$b$} (013);
  \path[->] (013)  edge [above,pos=.6] node {$b$} (012);
  \path[->] (012)  edge [bend left,left] node {$b$} (123);
  
  \path[->] (023)  edge [above] node {$a$} (13);
  \path[->] (13)   edge [bend left,left] node {$b$} (02);
  \path[->] (02)   edge [bend left=40,right] node {$b$} (13);
  \path[->] (13)   edge [bend left,right] node {$a$} (23);
  \path[->] (23)   edge [bend left,left] node {$a$} (13);
  \path[->] (23)   edge [above] node {$b$} (03);
  \path[->] (03)   edge [right] node {$b$} (01);
  \path[->] (01)   edge [above] node {$a,b$} (12);
  \path[->] (12)   edge [left] node {$b$} (23);
  \path[->] (03)   edge [above] node {$a$} (13);
  \path[->] (12)   edge [loop below] node {$a$} (12);
  
  \path[->] (02) edge [bend left,above] node {$a$} (1);
  
\end{tikzpicture}}%
\end{minipage}%
 %\caption{Automaton from Example~\ref{ex:I_vs_Gamma1} where $I(\TS{\mathcal A})$ 
% 	is a proper subgraph of $\Gamma_1(\TS{\mathcal A})$.}
\caption[Example power automaton]{Left: An automaton. Right: Its power automaton (empty set omitted).}
 \label{fig:thm_not_usable}
\end{figure}

\begin{proposition} \label{prelim:fact:cycle_states} \label{prelim:automata:power:P}
	Let $\semiautomat{A}{Q}{\Sigma}{\delta}$ with $n = |Q|$ and $u \in \Sigma^*$. 
	Then for $i \ge 0$ we have
	$\delta(Q, u^{i+1}) \subseteq \delta(Q, u^i)$
	and for $P = \delta(Q, u^{n-1})$
	the following is true:
	\begin{enumerate}
	\item $P = \delta(Q, u^{n - |P|})$,
		
	\item $\delta(P, u) = P$ ($u$ permutes~$P$) where~$P$
	is the largest set with this property, and $P = \delta(Q, u^i)$ for all $i \ge n - |P|$,
		
	\item $P \subseteq \delta(Q, u^i)$ for all $i \ge 0$.
	\end{enumerate}
\end{proposition}
\begin{proof}
	For the first claim, since $\delta(Q, u) \subseteq Q$
	and $\delta(Q, u^{i+1}) \subseteq \delta(Q, u^i)$
	implies 
	\[ \delta(\delta(Q, u^{i+1}), u) \subseteq \delta(\delta(Q, u^i), u),
	\]
	by induction we find that $\delta(Q, u^{i+1}) \subseteq \delta(Q, u^i)$
	for all $i \ge 0$.
	
	Next, by finiteness, there exists a minimal $j \ge 0$
	such that $\delta(Q, u^{j+1}) = \delta(Q, u^j)$.
	Then we have $\delta(Q, u^k) = \delta(Q, u^j)$ for all $k \ge j$
	and we get the following chain of proper inclusions
	\begin{equation}\label{prelim:automata:power:eqn:seq}
		\delta(Q, u^j) \subsetneq \delta(Q, u^{j-1}) \subsetneq\cdots \subsetneq \delta(Q, u) \subsetneq Q.
	\end{equation}
	This yields $\delta(Q, u^j) \subseteq \delta(Q, u^i)$ for every $i \ge 0$.
	Since $|\delta(Q, w)| \ge 1$ for all $w \in \Sigma^*$,
	the above chain of proper inclusions has length at most $n - |\delta(Q, u^j)|$, i.e., $j  \le n - |\delta(Q, u^j)|$.
	Hence we find
	\[
	P = \delta(Q, u^{n-1}) = \delta(Q, u^{n - |\delta(Q, u^j)|}) = \delta(Q, u^j)
	= \delta(Q, u^{n}) = \delta(P, u),
	\]
	and so $P = \delta(Q, u^{n - |P|})$,
	$P \subseteq \delta(Q, u^i)$ for all $i \ge 0$,
	and $P$ is permuted by~$u$. This also gives
	$P = \delta(Q, u^i)$ for $i \ge n - |P|$
	by induction.	
	Furthermore, if $\delta(S, u) = S$ for $S \subseteq Q$,
	then $S = \delta(S, u^i) \subseteq \delta(Q, u^{i})$ for all $i \ge 0$.
	In particular, we obtain $S \subseteq \delta(Q, u^{n-1}) = P$
	and so $P$ is the largest set permuted by~$u$. This shows all claims.
\end{proof}

The set~$P$ from \cref{prelim:fact:cycle_states}
is precisely the set of cycle states described in \cref{prelim:rem:intuition-functional-graph}.
By finiteness, after applying a word at most $n - 1$ times
to any state we end up at a cycle state. 

Note that since $\delta(Q, u) \subseteq Q$ for $u \in \Sigma^*$,
we have $\delta(Q, uv) \subseteq \delta(Q, v)$ for each $v \in \Sigma^*$.

\begin{corollary}  \label{prelim:automata:cor:power:P}
	Let $\semiautomat{A}{Q}{\Sigma}{\delta}$ with $n = |Q|$ and $u \in \Sigma^*$. 
	Then we have $\delta(Q, u^i) = \delta(Q, u^{\min\{n-1,i\}})$ for $i \ge 0$.
\end{corollary}
\begin{proof}
	 By \cref{prelim:automata:power:P}, the word~$u$
	 permutes the set $\delta(Q, u^{n-1})$.
	 Thus, for $i \ge n - 1$ we find $\delta(Q, u^i) = \delta(Q, u^{n-1})$, which gives the claim.
\end{proof}

In~\cite[Lem.~1~\&~2]{DBLP:journals/iandc/PribavkinaR11} it was also proved that $P$
is the largest set permuted by~$u$ in \cref{prelim:fact:cycle_states},
together with the fact that we have to apply~$u$ at most $n - |P|$
times to~$Q$ to obtain~$P$.

Note that for the properties collected in \cref{prelim:automata:power:P} to hold
it is essential that we start from the whole state set before applying powers of~$u$.

\begin{proposition} \label{prelim:word_in_PA} \label{prelim:automata:power:prop:S}  
Let $\semiautomat{A}{Q}{\Sigma}{\delta}$, $n = |Q|$, $u \in \Sigma^*$, $S \subseteq Q$, and $i \ge 0$.

\begin{enumerate}
\item \label{prelim:automata:power:prop:S:1} $\delta(S, u^i) = \delta(S, u^j)$ for some $0 \le j < 2^n - 1$.

\item \label{prelim:automata:power:prop:S:2} $|\delta(S, u^i)| = |\delta(S, u^j)|$ for some $0 \le j \le n - |\delta(S, u^i)|$. 

\item \label{prelim:automata:power:prop:S:3} $|\delta(S, u^i)| = |\delta(S, u^{\min\{n-1,i\}})|$. 
\end{enumerate}
\end{proposition}
\begin{proof}
\ref{prelim:automata:power:prop:S:1}: We can assume $S \ne \emptyset$.
 Let $j \ge 0$ be a minimal number
 such that $\delta(S, u^i) = \delta(S, u^j)$.
 Assume $j \ge 2^n - 1$.
 We have $\delta(S, u^k) \ne \emptyset$ for all $k \ge 0$. 
 Since there are at most $2^n - 1$ non-empty subsets, 
 and the sequence $\delta(S, u^0), \delta(S, u), \ldots, \delta(S, u^{2^n-1})$
 has length $2^n$, there exist $0 \le k < k' \le 2^n - 1$
 such that $\delta(S, u^{k}) = \delta(S, u^{k'})$.
 Then $\delta(S, u^{j - (k' - k)}) = \delta(S, u^j)$,
 contradicting the minimal choice of~$j$.
 So, we have $j < 2^n - 1$.

\ref{prelim:automata:power:prop:S:2}: We can assume $i > n - |\delta(S, u^i)|$.
 For the minimal $j \ge 0$ 
 with $\delta(Q, u^{j+1}) = \delta(Q, u^j)$, 
 we have a proper chain of inclusions as in \cref{prelim:automata:power:eqn:seq}.
 Thus, for $k \in \{0,1,\ldots,j\}$,
 we get $|\delta(Q, u^k)| \le n - k$.
 Since $\delta(S, u^j) \subseteq \delta(Q, u^j)$,
 we find $|\delta(S, u^j)| \le n - j$,
 and so $j \le n - |\delta(S, u^j)|$.
 Set $P = \delta(Q, u^{n-1})$.
 With \cref{prelim:automata:power:P} we get
 $\delta(Q, u^j) = P$ and that $u$ permutes~$P$,
 and so $u$ acts injectively on subsets of~$P$.
 Thus, from $\delta(S, u^j) \subseteq P$
 we obtain $|\delta(S, u^k)| = |\delta(S, u^j)|$
 for all $k \ge j$.
 In particular, since $i > j$,
 we find $|\delta(S, u^i)| = |\delta(S, u^j)|$.
  
 \ref{prelim:automata:power:prop:S:3}: We can assume $i \ge n$.
 	By Item~\ref{prelim:automata:power:prop:S:2} 
 	we have $|\delta(S, u^i)| = |\delta(S, u^j)|$
 	for some $0 \le j < n$.
 	With \cref{eqn:subset-mapped-to-smaller-set} 
 	we find
 	\begin{enumerate}
 	\item $
 	|\delta(S, u^{n-1})| = |\delta(\delta(S, u^j), u^{n - 1 - j})| \le |\delta(S, u^j)|
 	$
 	and
 	
 	\item $|\delta(S, u^i)| = |\delta(\delta(S, u^{n - 1}), u^{i - n + 1})| \le |\delta(S, u^{n - 1})|.
 	$
 	\end{enumerate}
 	Thus we obtain $|\delta(S, u^i)| = |\delta(S, u^{n-1})|$.
\end{proof}

\begin{corollary} 	\label{prelim:fact:unary_sync_word}
 Let $\semiautomat{A}{Q}{\Sigma}{\delta}$, $u \in \Sigma^*$, and $S \subseteq Q$.
 Then $|\delta(S, u^i)| = 1$ for some $i \ge 0$
 if and only if $|\delta(S, u^{n - 1})| = 1$ for $n = |Q|$.
\end{corollary}
\begin{proof}
  First, assume $|\delta(S, u^i)| = 1$ for some $i \ge 0$.
  Since $|\delta(S, u^{i})| = |\delta(S, u^{\min\{n-1, i\}})|$
  by \crefitem{prelim:automata:power:prop:S}{prelim:automata:power:prop:S:3} 
  and $|\delta(S, u^{n - 1})| \le |\delta(S, u^{\min\{n-1, i\}})|$, 
  we find $|\delta(S, u^{n - 1})| = 1$.
  The converse implication is obvious.
\end{proof}

\begin{remark}
To obtain \crefitem{prelim:automata:power:prop:S}{prelim:automata:power:prop:S:1},
we have basically used the fact that in any $n$-state automaton 
we enter a cycle when we read a word more than $n-1$ times.
Now, as $u$ permutes $P$ in \cref{prelim:fact:cycle_states},
the smallest $i \ge 0$ such that $\delta(q, u^i) = q$
for all $q \in P$
is asymptotically equal to $\exp((1+o(1))\sqrt{(|P|\ln{|P|})}$
by a result from Landau~\cite{landau03}.
This can be used to improve the bound from \cref{prelim:word_in_PA}
asymptotically. However, as 
the bound given is sufficient for our applications,
we do not bother to state a more precise bound.
\end{remark}